\cotefile{27}
\scanpage{15}

$1^\circ)$ On peut prendre pour $n$ le degré de \unsure{tte|toute} extension
finie $k'$ de $k$ ayant les propriétés suivantes

a) $G_{k'\,\mathrm{réd}}$ est lisse sur $k'$ i.e. séparable $/k'$

b) \deleted{$G_{k'\,\mathrm{réd}}$ admet un Borel}
les pts de $G/G^0$ sont \deleted{définis}\inserted{rationnels} \uline{sur $k'$}.

c) \inserted{``le'' radical $R'$} \unsure{de} $G^0_{k'\,\mathrm{réd}}$ \deleted{est} \inserted{est défini sur $k'$, et splitté}
\deleted{admet un Borel $B'$ qui est splitté} sur $k'$ [suite de composition à
quotients successifs $G_a$ ou $G_m$]

d) \inserted{le \unsure{gpe rédt|groupe rédt|hyp} $G^0_{k'\,\mathrm{réd}}/R'$ est \deleted{un groupe \unsure{réductif|$k$ simple}} déployé}
\deleted{\unsure{de|le} $T'$ est un tore maximal de $B'$}
\deleted{[nécessairement splitté] et $N(T')$, $C(T')$}
\deleted{\unsure{leur|le|son} normalisateur resp\textsuperscript{t} centralisateur, dans}

\deleted{Par suite,}

\deleted{$1^\circ)$ Si $G$ est unipotent \inserted{connexe}, on peut prendre
$n$ une puissance de la caractéristique}

\deleted{$2^\circ)$ Si $G$ est \deleted{réductif} $\frac12$ simple, alors
($\operatorname{Hom}(G,G_m)$ étant nul, après \deleted{ext} \unsure{ext.} du
corps de base] $\operatorname{Pic}$}
\deleted{$G$ est lisse}

$2^\circ)$ Supposons que $\operatorname{Hom}_{\unsure{\bar k\text{-}\mathrm{gr}|\bar k\text{-}\mathrm{gp}}}(G_{\bar k},G_{m\,\bar k})=0$,
et que $G$ soit \inserted{lisse et} connexe. Soit $\widetilde k$ la
clôture \inserted{algébrique} \uuline{séparable} de $k$. Alors $\operatorname{Pic}(G)\simeq\operatorname{Pic}(G_{\widetilde k})$

\deleted{En particulier, si $\operatorname{Hom}(G_{\unsure{k|\bar k}},G_m)$}

Si \inserted{on a} $\underline{\operatorname{Hom}}_{\unsure{\widetilde k\text{-}\mathrm{gr}|\widetilde k\text{-}\mathrm{gp}}}(G_{\widetilde k},G_{m\,k})=0$ [ce qui est
vrai si en plus de la \unsure{relation|condition} \unsure{précédte|précédente} sur $\bar k$,
$G$ est de \uline{rang unipotent} \uline{nul}] on a $=\operatorname{Pic}(G)\simeq\operatorname{Pic}(G_{\widetilde k})$

% archive sheet number: 14
\scanpage{16}

$3^\circ)$ Si $G$ est \uline{unipotent}, alors on \deleted{\unsure{que|a|sait que}} $\operatorname{Pic}(G)$
est annulé par une puissance de la car. $p$.

Si $G$ est de rang unipotent nul et $\operatorname{Hom}_{\unsure{\bar k\text{-}\mathrm{gr}|\bar k\text{-}\mathrm{gp}}}(G_{\bar k},G_{m\,\bar k})
=0$, alors $\operatorname{Pic}(G)$ \deleted{\unsure{est}} $\unsure{=|\simeq}$ \deleted{fini} $\operatorname{Pic}(G_{\bar k})={}$ \uline{\unsure{gpe|groupe|grpe}} \uline{fini}

% archive sheet number: 15
\scanpage{18}

\[
\deleted{(\mathrm{vii})\overset{?}{\Longrightarrow}(\mathrm{vi})}
\]

\margin{?}
\diagram{\resizebox{\linewidth}{!}{\begin{tikzcd}[ampersand replacement=\&,column sep=0.7em,row sep=0.4ex,cells={nodes={font=\small}}]
\deleted{(\mathrm{vii})\not\Longrightarrow(\mathrm{vi})}
 \arrow[dash,from=1-1,to=7-3,start anchor=north east,end anchor=south west]
 \arrow[dash,from=1-2,to=7-1,start anchor=north west,end anchor=south west]
 \arrow[dash,from=1-3,to=7-2,start anchor=north east,end anchor=south east]
 \arrow[dash,from=1-3,to=7-4,start anchor=north west,end anchor=south east]
 \arrow[dash,from=2-4,to=7-3,start anchor=north east,end anchor=south east]
 \& \deleted{\text{\={m} si }\psi\ \text{\uline{birationnel}}} \& {} \& \deleted{\mathrm{Sl}(2,\unsure{k|K})\ \text{\uline{car. 2.}}}\\
\deleted{(\mathrm{vi})\not\Longrightarrow(\mathrm{v})} \& {} \& {} \& {}\\
\deleted{(\unsure{\mathrm{x}|\mathrm{v}})\not\Longrightarrow(\mathrm{i})} \& {} \& {} \& {}\\
\deleted{(\mathrm{vi})\not\Longrightarrow(\mathrm{x})\quad\text{O.K.}} \& {} \& {} \& {}\\
{} \& |[draw,rounded corners=7pt]| \deleted{\inserted{\text{à fortiori }(\mathrm{vi})\not\Longrightarrow(\mathrm{i}),}}
 \& |[draw,rounded corners=7pt]| \deleted{\inserted{(\mathrm{x})\not\Longrightarrow(\mathrm{xiii})}} \& {}\\
\deleted{(\mathrm{x})\not\Longrightarrow(\mathrm{i})} \& \deleted{\text{Cas semi-simple \textendash\textendash\textendash\textendash\textendash}} \& {} \& {}\\
\deleted{(\mathrm{i})\not\Longrightarrow(\mathrm{x})} \& |[draw,rounded corners=7pt]| \deleted{\text{à fortiori }(\unsure{\mathrm{v}|\mathrm{vi}})\not\Longrightarrow(\mathrm{v}),}
 \& |[draw,rounded corners=7pt]| \deleted{(\mathrm{vi})\not\Longrightarrow(\mathrm{x}),}
 \& \deleted{\unsure{\mathrm{GP}|\mathrm{PGL}}(2,\bar k),\ \text{\uline{car. 2}}}\\
{} \& {} \& |[draw,rounded corners=7pt]| \deleted{(\mathrm{i})\Longrightarrow(\mathrm{xiii})} \& {}
\end{tikzcd}}}

\noindent\rule{\linewidth}{0.4pt}

\margin{!}

{\small
\[
\left\{\begin{aligned}
&(\mathrm{vii})\not\Longrightarrow(\mathrm{vi})\quad\text{\={m} si $\psi$ birationnel :}\quad\mathrm{Sl}(2,k),\ \text{car. }2)\\[3pt]
&(\mathrm{x})\not\Longrightarrow(\mathrm{v})\quad[\mathrm{x}\not\Longrightarrow(\mathrm{i}),\
  (\mathrm{x})\not\Longrightarrow(\mathrm{xiii}),\ (\mathrm{vi})\Longrightarrow(\mathrm{v})]\\
&\hspace{4em}\text{Cas d'une \unsure{représn|repr\textsuperscript{on}|système} de $G$ semi-simple}\\[-2pt]
&\hspace{18em}\text{\textendash\textendash\textendash\textendash}\\
&(\mathrm{i})\not\Longrightarrow(\mathrm{x})\quad[(\mathrm{v})\not\Longrightarrow(\mathrm{x}),\
(\mathrm{vi})\not\Longrightarrow(\mathrm{x}),\ (\mathrm{i})\not\Longrightarrow(\mathrm{xiii})]\\
&\hspace{7em}\mathrm{GP}(2,k),\ \text{car. }2
\end{aligned}\right.
\]}

\margin{\deleted{$G=\unsure{\mathrm{P}\ell|\mathrm{Pl}|\mathrm{PGl}}(\unsure{1|2},k)$}\\
\deleted{$\dim N=1$}\\
\deleted{$\dim X:3-1+2=4$}}

\[
\boxed{(\mathrm{v})\not\Longrightarrow(\mathrm{i})\quad???}\qquad \mathrm{x}+\mathrm{i}
\]

\[
\tikz[baseline=(a.base)]{\node[draw,rounded corners=15pt,inner sep=5pt] (a) {$\begin{aligned}
(\mathrm{vii})&\not\Rightarrow(\mathrm{vi})\\
(\mathrm{x})&\not\Rightarrow(\mathrm{v})\\
(\mathrm{v})&\not\Rightarrow(\mathrm{i})
\end{aligned}$};}
\qquad
\tikz[baseline=(b.base)]{\node[draw,rounded corners=12pt,inner sep=5pt] (b) {$\mathfrak h=M$};}
\]

\[
(\unsure{\mathrm{xiii}|\mathrm{xii}})=(\mathrm{x})+(\mathrm{vii})
\]

\[
\boxed{(\mathrm{xiii})=(\mathrm{xii})+(\mathrm{v})\deleted{\unsure{\mathrm{ii}|?)}}}
\]

% archive sheet number: 17
\scanpage{19}

\deleted{Si $\mathfrak g$ \unsure{satisfait|satisfaisait}}

\uline{Lemme}\qquad $\mathfrak h$ alg. de Lie nilpotente \inserted{sur corps \uline{infini} $k$} opérant sur
vectoriel $V$. Supposons que pour $\forall\,x\in\mathfrak h$, $\rho(x)$
soit non injectif, Alors $\exists\,v\in V,\ v\ne0$ tel que
$\rho(\mathfrak h)\,v=0$

En effet, $V=\coprod V_i$, $V_i$ \deleted{propre}\inserted{primaire} pour la
fonction poids $\lambda_i$ sur $\mathfrak h$. L'hypothèse signifie
que pour $\forall\,\unsure{\dot x|x}\in\mathfrak h$, $\exists\,i$ \unsure{tels|tel} que $\lambda_i(x)=0$,
donc $V=\bigcup V(\lambda_i)$. Par le principe du
\unsure{prolongemt|prolongement} des identités, on en conclut que
$\exists\,i$ tel que $\lambda_i=0$. Donc, $\mathfrak h$ opère sur \deleted{$V_i$}
\deleted{\unsure{avec seul}} par opérations nilpotentes. Donc
par Engel, $\exists\,v\in\mathfrak h,\ v\ne0$, tel que $\mathfrak h.v=0$, O.K.

\[
\unsure{f|P}_i(X,t)=\sum\unsure{C_i|G_i|c_i}(x)t^{n-i}
\]

\[
\underbrace{t^{n}+\unsure{C_1(x)|G_1(x)|c_1(x)}t^{n-1}\ \unsure{\cdots|{-}|\ldots}\ \unsure{+|-|\pm}C_n}_{[t-\lambda(x)]^{\unsure{n|m}}}=0
\]

\[
\lambda(x)^n=C_n
\]

% archive sheet number: 18
\scanpage{20}

\[
\mathscr L\subset\mathfrak g
\]

\[
\mathscr L+k\mathfrak g_{-\alpha}\qquad\sum_{\beta\ne\alpha}\lambda_\beta X_{-\beta}
\]

\[
\deleted{\unsure{X_\beta|\lambda_\beta}}\qquad\bar\beta=0
\]

\[
\left[X_{-\alpha},\sum_{\beta\ne\alpha}\lambda_\beta X_{-\beta}\right]
=\sum\lambda_\beta\left[X_{-\alpha},X_{-\beta}\right]
\]

\[
[X_\alpha,X_\beta]
\]

Deux racines \qquad $X_1,X_2,\unsure{X_3|N_3}$

\[
\begin{array}{cc}
\begin{pmatrix}\cdot&\\ &\cdot\end{pmatrix}
&
\begin{pmatrix}\cdot&\cdot&\cdot\\ &\cdot&\cdot\\ &&\cdot\end{pmatrix}\\[5pt]
\mathscr L & \deleted{\unsure{\mathfrak b|\mathscr L}}
\end{array}
\]

\diagram{\begin{tikzcd}[ampersand replacement=\&,column sep=small,row sep=small]
(\mathrm{xiv})\arrow[r,Rightarrow]\arrow[dr,Rightarrow] \&(\mathrm{i})\arrow[r,Rightarrow] \&(\mathrm{iii})\arrow[r,Rightarrow] \&(\mathrm{v})\\
\&\underline{(\mathrm{x}\deleted{\unsure{\mathrm{ii}|\mathrm{i}}})}\arrow[r,Rightarrow] \&(\mathrm{xii})
\end{tikzcd}}

\noindent\begin{tikzpicture}
\node[inner sep=5pt] (casebox) {\begin{minipage}{0.91\linewidth}

\[
\left[\begin{aligned}
&\text{Une \uline{seule} racine } \alpha>0\quad\text{telle que } \bar\alpha=0\\
&\text{Une racine } \beta>0\quad\text{telle que } [X_\alpha\unsure{,|:}X_\beta]\ne0
\end{aligned}\right.
\]

Alors $\mathscr L+k\mathfrak g_{-\beta}$ est son propre normalisateur,
\uline{satisfait} au th. de densité
mais ne contient \unsure{pas|que} de \unsure{sous-|non-}algèbre de Cartan

\end{minipage}};
\draw[rounded corners=4pt] (casebox.north west) -- (casebox.north east) -- (casebox.south east) -- ([xshift=28mm]casebox.south west) -- ([xshift=28mm,yshift=2mm]casebox.south west) -- ([yshift=2mm]casebox.south west) -- cycle;
\end{tikzpicture}

% archive sheet number: 19
\scanpage{21}

les sous-a\deleted{\unsure{g|gg}}lgèbres égales à leur propre normalisateur).

Démonstration. a) Si \hspace{1.5em}$\operatorname{Nil}(a\deleted{\unsure{,|.}})\quad)$, comme cette dernière est nil\unsure{p}potente,
\hspace{1em}l'est. Inversement, si \hspace{1.1em} est nilpotent, elle est contenue dans
le nil-espace de son élément $a$, donc \hspace{2em}$\operatorname{Nil}(a\deleted{,})\quad)$.

b) Si \hspace{2em}$\operatorname{Nil}(a\deleted{,})\quad)$, comme $\operatorname{ad}(a)_{\quad/\operatorname{Nil}(a\deleted{,})}$ est injectif, il en
est à fortiori de $\operatorname{ad}(a)_{\quad/\quad}$, donc \hspace{1.5em} est son propre normalisateur
connexe. \deleted{\unsure{En|In}}Inversement, si \hspace{1em} est son propre normalisateur connexe, alors

% archive sheet number: 20
\scanpage{22}

\uline{\unsure{Remarques|Remarque}}\quad 1) \={M} si $\psi$ est \deleted{\unsure{gén.} étale}\inserted{\uline{birationnel}}, \uline{et} $\mathfrak h$ provient
d'un sous-groupe $H$ de $G$ lisse sur $k$, \deleted{\unsure{la même|aussi}}
la condition (i) n'est pas néc. vérifiée.

\uline{Exemple}\quad $G=\mathrm{Sl}(2,k)$\quad $k$ car. 2, Alors
$\mathfrak g$ est \uline{nilpotent} donc \deleted{\unsure{les}}la \inserted{seule} sous-algèbre de
Cartan \deleted{\unsure{de sont|de son}} est $\mathfrak g$ elle \={m}. Soit
$B$ Borel, alors \deleted{$\mathscr L$} $\mathfrak h=\operatorname{Lie}(B)$, on sait
que $\bigcup_{g\in G(k)}\operatorname{ad}(g)\mathscr L=\mathfrak g$ \deleted{\unsure{, et}} \unsure{(cf)|(sic)}, \deleted{\unsure{et|donc}}
\deleted{\unsure{est}} donc $\psi$ dominant. D'ailleurs
$M=\mathfrak g\quad B\subset N\ne G$, donc $\dim N=\dim B=\operatorname{rang}\mathfrak h$,
donc $\dim X=\dim G$, \deleted{D'ailleurs on a} \inserted{donc $\psi$ gén. quasi fini.}

\deleted{aussi pour $\underline{\forall}\ a\in\mathfrak h$, $M=M_a=\mathfrak g$\quad (car $M\subset M_a$)}

D'ailleurs, une analyse plus poussée montre
\deleted{donc $\varphi$ \uline{lisse} en tout point \unsure{gén.} de \unsure{$X$}}
\deleted{$\psi$ étale en tout point}

\diagram{\begin{tikzcd}[ampersand replacement=\&,row sep=small,column sep=small]
\& \deleted{N=M_{a}} \& \\
\deleted{\unsure{G/N|G_N}} \& \deleted{\unsure{M=M_a|M=\mathfrak g}} \& \deleted{\begin{pmatrix}a&b\\0&-a\end{pmatrix}}\\
\deleted{\unsure{G|\mathscr L}\quad B\subset G} \& \&
\end{tikzcd}}

ici que $\psi$ \uline{est}
\deleted{\unsure{quasi-ét.}} \uline{birationnel}
\deleted{\unsure{un système de B et B$^-$|$B\cap B$}}

\reversemarginpar
\margin{\setlength{\fboxsep}{1pt}\fbox{\parbox{0.94\marginparwidth}{\raggedright\scriptsize\uline{N.B.}\quad $\varphi$ lisse en $(\dot e,a)$\\
$\Updownarrow$\\
$\operatorname{ad}(a)_{\mathfrak g/\mathfrak h}$ injectif\\
i.e. $[x,a]\in\mathfrak h\Rightarrow x\in\mathfrak h$\\
$M_a=\mathfrak h$\\
$\Downarrow$\\
$\mathfrak h=M$}}}
\normalmarginpar

2) \={M} si $\psi$ \deleted{birationnel}\inserted{dominant} et $\mathfrak h=M$
la condition (\unsure{x|i|xi}) n'est pas
néc. vérifiée. Prenons
$G=\mathrm{Pl}(2,k)$, $k$ car 2, $\mathfrak g$ a une
base $X,Y,H$ avec
$[X,Y]=2H=0,\quad[H,X]=X,\quad[H,Y]=Y$.

% archive sheet number: 21
\scanpage{23}

\margin{\tikz[baseline=(r.base)]{\node[inner sep=1pt] (r) {$\mathfrak h\subset M\subset M_a$};\draw (r.north west)--(r.south west)--(r.south east);}}

% Three diagonal strokes run through the draft. The dots at the right are ink.
{\scriptsize
\diagram{\begin{tikzcd}[ampersand replacement=\&,row sep=2pt,column sep=2pt,cells={anchor=west}]
{}\arrow[from=1-1,to=6-1,draw=none,"{\left\{\vphantom{\rule{0pt}{4.4em}}\right.}"{left,inner sep=1pt}]
\arrow[from=9-1,to=12-1,draw=none,"{\left\{\vphantom{\rule{0pt}{3.7em}}\right.}"{left,inner sep=1pt}]
\arrow[dash,from=1-2,to=12-2,start anchor=north west,end anchor=south]
\arrow[dash,from=1-2,to=10-2,start anchor=north,end anchor=south east] \&\deleted{(\mathrm{i})\ ,\ \mathfrak h\text{ contient une sous-algèbre de Cartan}}\\
{} \&\deleted{(\mathrm{ii})\quad\mathfrak h=M,\quad\mathfrak h\text{ contient \unsure{élémt|élément|des éléments} régulier de }\mathfrak g}\\
{}\arrow[ddd,Rightarrow] \&\deleted{(\mathrm{iii})\quad\exists\,a\in\mathfrak h\quad\text{t.q. }\operatorname{ad}_{\mathfrak g/\mathfrak h}(a)\text{ injectif}\quad
(\text{i.e. tel que }\deleted{\unsure{M_a|M}}\mathfrak h=M_a)}\\
{} \&\deleted{(\mathrm{iv})\quad\psi:X\unsure{\mathbin{-}|\longrightarrow}W(\mathfrak g)\quad\text{dominant, et }\mathfrak h=M\qquad\bullet}\\
{} \&\deleted{(\mathrm{v})\quad\psi:X\mathbin{-}W(\mathfrak g)\quad\text{gén. lisse}}\\
{} \&\deleted{(\mathrm{vi})\quad\varphi:X\times W(\mathfrak h)\longrightarrow W(\mathfrak g)\quad\text{gén\textsuperscript{t} lisse}}\\
{} \&\deleted{(\mathrm{vii})\quad\psi:X\longrightarrow W(\mathfrak g)\quad\text{dominant}\qquad\bullet}\\
{}\arrow[u,Rightarrow] \&\deleted{(\mathrm{viii})\quad\psi:X\mathbin{-}W(\mathfrak g)\quad\text{dominant et quasi-fini}\qquad\bullet}\\
{}\arrow[u,Rightarrow] \&\deleted{(\mathrm{ix})\quad\psi\quad\text{gén. \unsure{qu.|quasi} fini et }\mathfrak h=M}\\
{}\arrow[uuuuuuu,Rightarrow,bend left=25] \&\deleted{\deleted{(\unsure{\mathrm{viii}|\mathrm{xiii}})}(\mathrm{x})\qquad\deleted{\psi\text{ dominant}}\quad\psi:X\longrightarrow W(\mathfrak g)\quad\text{dominant, \uline{et}}\quad\bullet}\\
{} \&\deleted{\left[\begin{aligned}&\mathfrak h\text{ est l'alg. de Lie d'un sous-groupe lisse }H\\&\text{de }\mathfrak g\end{aligned}\right.}\\
{} \&\deleted{\deleted{(\unsure{\mathrm{x}|\mathrm{ix}|\mathrm{xii}})}(\mathrm{xi})\quad\psi:X\mathbin{-}W(\mathfrak g)\quad\text{dominant, et }N\text{ lisse.}\qquad\bullet}
\end{tikzcd}}}

\deleted{$\mathfrak h$ \deleted{\unsure{plus}} minimale parmi les $k$. algèbres de Lie qui sont}

\deleted{a) soit $={}$ normalisateur, et \unsure{satisfont|satisfait|satisfasse} \unsure{au|aux} th. de densité}

\deleted{b) soit $={}$ normalisateur, et satisfont aux th. des \unsure{dérivations|désintégrations}}

\deleted{c) soit définies par $H$ lisse, et \rule{5em}{0.3pt} densité}

\deleted{d) soit telles que $\exists\,a\in\mathfrak h\quad\operatorname{ad}(a)_{\mathfrak g/\mathfrak h}$ injectif}

\deleted{e) soit telles que}

% archive sheet number: 22
\scanpage{24}

Ici, \inserted{la sous-algèbre de Cartan est $\unsure{D|\mathfrak d|\delta}=\mathfrak t=kH$.}
\unsure{toute|tout} sous espace \deleted{de la forme} vectoriel de
$\mathfrak g$ contenant $\mathfrak t=kH$ \deleted{\unsure{la réponse du gpe à}} \inserted{est une sous-algèbre.}
\deleted{\unsure{faire varier le type}, $=$ Cartan}

Mais prenons $\mathfrak h=kH+k(\lambda X+\mu Y)$, si \deleted{$\mathfrak h$ était}
\deleted{\unsure{réalisé}} le normalisateur \inserted{$N$} de $\mathfrak h$ étant de dim 2,
\deleted{il contiendrait $T'$} \inserted{donc $\mathfrak h=\operatorname{Lie}(H)$ i.e. lisse} \deleted{donc} alors
\deleted{$H$ contiendrait un tore maximal $T'$, \unsure{donc}}
\deleted{\unsure{un conjugué de}} serait résoluble, donc un
\deleted{\unsure{visiblement facile}}\quad \deleted{le Borel}
sous-groupe de Borel, or \deleted{les seuls sous-groupes}
($T$ \unsure{étant|est} un C-sous-groupe) \deleted{\unsure{donc}} $\mathfrak h\supset\mathfrak t$

\deleted{le Borel-type $B$, tel que}
\deleted{$\mathscr L=kH+kX$, est tel que $B=\operatorname{Norm}\mathscr L$,}
\deleted{donc est \unsure{défini}}

impliquerait $H\supset T$,
\deleted{donc} or il n'y a que deux Borels
contenant $T$, \deleted{\unsure{si}} donc si $\lambda,\mu\ne0$, on
voit que $\unsure{\mathfrak h|\mathfrak g}$ n'est pas Algèbre de Lie d'un sous-groupe lisse de $G$ .....

3)

% archive sheet number: 23
\scanpage{25}

Exemple d'un groupe lisse résoluble \inserted{connexe} $G$ sur $\mathbb F_p$, \inserted{de rang unipotent nul \deleted{\unsure{réductif|scindé} \illegible{1}} \unsure{linéaire|trivial},} dont l'algèbre
de Lie n'a pas d'élément régulier.

Soit $(\alpha_i)_{i\in I}$ une famille \inserted{finie} de formes \inserted{lin. $\ne0$} sur $\deleted{\unsure{\mathbb F_p^2|\mathbb F_p}}\inserted{\mathfrak t=\mathbb F_p^2}$ telles que
$\bigcup\operatorname{Ker}\alpha_i=\deleted{\unsure{\mathbb F_p|\mathbb F_p^2}(\mathbb F_p^2)}\mathfrak t\deleted{\mathbb F_p^2}$. (On peut s'en tirer
avec $p+1$ formes non nulles, et $\bigcap\operatorname{Ker}\alpha_i=0$). \deleted{Soit}
Les $\alpha_i$ sont les caractères infinitésimaux associés
à des \unsure{homs|hom|hom.}

\[
\chi_i:T=G_m^2\longrightarrow G_m
\]

qu'on utilise pour former le produit semi-direct

\[
G=T.U\deleted{G_a^I}\qquad U=G_a^I
\]

avec $T$ opérant sur $G_a^I$ par les $\chi_i$.
Alors $T$ opère sur $\mathfrak g/\mathfrak t$ par les caractères
$\chi_i$, ce qui montre que $T$ est un Cartan
i.e. le rang unipotent de $G$ est nul. Le
centre de $G$ est donc égal au centre réductif, donc
égal au noyau de $T\longrightarrow\operatorname{Gl}(\mathfrak g)$\,; qui est
nul, \deleted{\unsure{donc par suite}} Enfin, un \deleted{\unsure{par suite}} élément
\deleted{régulier} de $\mathfrak t$ \deleted{serait dans \unsure{le|la}} étant dans
$\bigcup\operatorname{Ker}\alpha_i$, ne peut être \uline{régulier}. Or \uline{$\unsure{\forall|\text{tt}|\text{tout}}$}
\uline{tore maximal} \uline{de $G$} est \uline{conjugué de} $T$
\deleted{\unsure{[comme|comme}} \inserted{(par un élément unique de $U(\mathbb F_p)$]} donc
\deleted{tout} $\mathfrak g$ n'a pas d'él\textsuperscript{mt} régulier [car
il serait conjugué d'un élément de $\mathfrak t$].

% slanted
\margin{\unsure{Ici|Donc}, $T$ est \deleted{\unsure{égal|éplt}} un sous-groupe de type $(C)$, i.e. $\mathfrak t$ est une sous-algèbre de Cartan}

% archive sheet number: 24
\scanpage{26}

\deleted{Soit}
Cas d'un schéma en groupes \inserted{séparé} étale de t.f. sur $S$
dont le centre n'est pas représentable
($S={}$spectre d'un anneau de valuation
discret $A$)

On part d'un groupe ensembliste $E$,
de centre trivial, et d'un sous-groupe commutatif $\deleted{\unsure{E|H}}F$ non \deleted{\unsure{nul}} \inserted{trivial} (p.ex. le groupe engendré
par un élément $g\ne e$ de $E$). Soit
$\deleted{\unsure{F|H|\mathcal F}}G$ le sous-groupe ouvert de $E_S$ obtenu
en enlevant de la fibre spéciale
$E_k$ de $E_S$ le fermé $E_k-\deleted{\unsure{E|H}}F_k$. Donc

\[
G_0=F_k,\qquad G_1=E_K
\]

Alors le foncteur $\operatorname{Cent}(G)$ n'est pas
représentable, car il ne commute
pas à la limite projective adique
des $A_n=A/\mathfrak m^{n+1}$, car on trouve

\[
\begin{aligned}
Z(\hat A)&=\{e\}\\
\varprojlim Z(A_n)&=F
\end{aligned}
\]

\uline{N.B.} Si $G$ est \inserted{plat} de présentation finie sur $S$, à fibres connexes,
\unsure{alors} $Z$ est représentable ......

% archive sheet number: 25
